% TOP_PROBLEM: {"id": 20003326, "title": "Arithmetic QUE through Hecke packets", "category_id": 20, "category": {"id": 20, "name": "miscellaneous", "display_name": "Miscellaneous", "description": "Problems whose source classification does not fit the main mathematical categories.", "slug": "miscellaneous", "order_index": 20, "created_at": "2026-07-31T15:26:25.671Z"}, "status": "partially_solved"}
% FIRST_POSTED: null
% !TeX program = lualatex
% Compile twice: lualatex 20003326.tex
% All references and prose are embedded; no BibTeX database or external inputs.
\documentclass[11pt,a4paper]{article}
\usepackage[margin=24mm,headheight=14pt]{geometry}
\usepackage{fontspec}
\setmainfont{Latin Modern Roman}
\setsansfont{Latin Modern Sans}
\setmonofont{Latin Modern Mono}
\usepackage{amsmath,amssymb,mathtools}
\usepackage{unicode-math}
\setmathfont{Latin Modern Math}
\usepackage{microtype}
\usepackage{enumitem,xurl,needspace}
\usepackage[dvipsnames]{xcolor}
\usepackage{hyperref}
\hypersetup{unicode=true,colorlinks=true,linkcolor=MidnightBlue,urlcolor=MidnightBlue,
 pdftitle={Research attempt: Problem 20003326},pdfauthor={Research notebook},bookmarksnumbered=true}
\usepackage{bookmark}
\usepackage{tocloft}
\setlength{\cftsecnumwidth}{3em}
\cftsetpnumwidth{2.5em}
\cftsetrmarg{4em}
\usepackage{titlesec}
\titleformat{\section}{\Large\bfseries\raggedright}{\thesection}{0.7em}{}
\usepackage{fancyhdr}
\pagestyle{fancy}\fancyhf{}
\fancyhead[L]{\small\sffamily Open Problems Math | Research notebook}
\fancyhead[R]{\small Problem 20003326 | 27 September 2026}
\fancyfoot[C]{\thepage}
\setlength{\parindent}{0pt}
\setlength{\parskip}{3pt plus 1pt}
\setlength{\emergencystretch}{3em}
\setcounter{tocdepth}{1}
\setcounter{secnumdepth}{1}
\setlist[itemize]{leftmargin=3em,itemsep=5pt,topsep=4pt,parsep=0pt}
\allowdisplaybreaks
\newcommand{\N}{\mathbb{N}}
\newcommand{\Z}{\mathbb{Z}}
\newcommand{\Q}{\mathbb{Q}}
\newcommand{\R}{\mathbb{R}}
\newcommand{\C}{\mathbb{C}}
\newcommand{\F}{\mathbb{F}}
\newcommand{\A}{\mathbb{A}}
\newcommand{\PP}{\mathbb P}
\newcommand{\E}{\mathbb{E}}
\newcommand{\eps}{\varepsilon}
\newcommand{\dd}{\,\mathrm d}
\newcommand{\sref}[2]{\hyperlink{source:#1:#2}{[S#2]}}
\newcommand{\eref}[2]{\hyperlink{extra:#1:#2}{[E#2]}}
% Turn the next switch off for a shorter reading copy. The quotations preserve
% every exactTarget field, including qualifications omitted from a short summary.
\newif\ifshowcatalogue
\showcataloguetrue
\newcommand{\cataloguescope}[1]{\ifshowcatalogue
 \paragraph{Exact scope in the supplied catalogue.}
 \begin{quote}\small #1\end{quote}\fi}
\usepackage{etoolbox}
\AtBeginEnvironment{itemize}{\small}
\numberwithin{equation}{section}
% Standalone export. Original research content preserved.
\begin{document}
\setcounter{section}{0}
% ================================================================
% problemId: 20003326
% Canonical problem ID: 20003326
% exactTarget SHA-256: d7734e1737549aeb99572740573dcafae6789c785447ce30c34087fcb92fcc71
\clearpage
\section*{\texorpdfstring{Arithmetic QUE through Hecke packets}{Arithmetic QUE through Hecke packets}}\label{20003326}
\subsection{Definitions and mathematical statement}
An arithmetic locally symmetric space has the form $\Gamma\backslash G/K_0$, with $\Gamma$ arithmetic and $K_0$ maximal compact. In the finite-volume source setting, take normalized joint Laplace/Hecke eigenfunctions $u_j$ with spectral parameters tending to high energy. Their probability densities are $|u_j|^2\,d\operatorname{vol}$. The stated conclusion is
\[
 |u_j|^2\,d\operatorname{vol}\ \rightharpoonup\
 \frac{d\operatorname{vol}}{\operatorname{vol}(\Gamma\backslash G/K_0)}.
\]
The phrase ``arithmetic quantum limit'' includes the specified joint-eigenfunction and spectral regime. The catalogue does not delimit all higher-rank or noncompact variants; those require source-specific hypotheses, including exclusion of escape of mass. It is not a blanket statement for arbitrary eigenfunctions on arbitrary arithmetic quotients.
\par\smallskip\textit{Formulation and concept references:} \sref{20003326}{1}.
\cataloguescope{Prove the arithmetic quantum unique ergodicity conjecture: for arithmetic locally symmetric spaces, the normalized Riemannian volume is the only arithmetic quantum limit of Hecke eigenfunctions.}
\subsection{Short English statement}
Do arithmetic high-energy joint eigenfunctions have only the uniform-volume distribution as a quantum limit?
% BEGIN RESEARCH ADDITION 137
\subsection{Research attempt}

\paragraph{Editorial note: a necessary scope distinction.}
Section 6.1, Conjecture 39 of the supplied AIM list was inspected, including
its definition of arithmetic quantum limits. Its brief formulation does
not expressly impose irreducibility of the lattice or growth of energy in
every factor of a reducible quotient. With arbitrary arithmetic products
and only total high energy, the conclusion fails by the construction below.
However, the catalogue explicitly retains the source's spectral regime.
Accordingly this is treated as a scope obstruction, \emph{not} a claimed
counterexample to an intended irreducible or nondegenerate version of
arithmetic QUE. The original wording is preserved. \sref{20003326}{1}

\paragraph{Current frontier.}
Shem-Tov--Silberman's recent theorem proves arithmetic QUE for congruence
quotients arising from norm-one quaternion groups over a number field,
including products of hyperbolic two- and three-spaces. Their Section 2.1
uses one number field in constructing the lattice, and Section 3 explicitly
uses its irreducibility. A Cartesian product of two independent arithmetic
surfaces is not that irreducible lattice setting. Their theorem therefore
does not eliminate the scope issue just identified. \eref{20003326}{1}
The broader measure-rigidity and mass-escape questions remain distinct
parts of the formulation source's discussion. \sref{20003326}{1}

\paragraph{Attack route.}
Test the universal statement under products before trying to invoke measure
rigidity. A fixed-energy factor can retain a nonuniform density while
another supplies all the growing energy. The proposed repair is factorwise
equidistribution; we prove exactly when that repairs a pure-tensor sequence.
On an irreducible lattice this factorization is absent, so the argument
cannot replace an arithmetic rigidity theorem.

\paragraph{Attempt: a product obstruction.}
Let $X_1,X_2$ be compact connected arithmetic hyperbolic surfaces from
indefinite quaternion division algebras and torsion-free congruence groups,
with probability volumes $m_1,m_2$. Use the commuting self-adjoint Hecke
family and the Laplacian; their eigenspaces have a simultaneous real
orthonormal eigenbasis. This compact arithmetic setup is included in the
source framework. \sref{20003326}{1}\eref{20003326}{1}
Choose normalized real joint eigenfunctions $\phi_j$ on $X_1$ with
$\lambda_j\to\infty$, and a fixed real nonconstant normalized joint
eigenfunction $\psi$ on $X_2$, with eigenvalue $\mu>0$. On the product set
\[
 u_j(x,y)=\phi_j(x)\psi(y).
\]
The product group and arithmetic lattice give an arithmetic locally
symmetric quotient. The product Hecke family and invariant differential
operators act jointly diagonally on $u_j$, and
\[
 \|u_j\|_{L^2(m_1m_2)}=1,
 \qquad -\Delta u_j=(\lambda_j+\mu)u_j,
 \qquad \lambda_j+\mu\longrightarrow\infty.
\]
These form a subsequence of a complete joint product eigenbasis, not merely
arbitrary Laplace eigenfunctions without Hecke symmetry.

The continuous test $a(x,y)=\psi(y)^2-1$ has volume mean zero, but
\begin{equation}\label{20003326:obstruction}
 \int a|u_j|^2\,dm_1dm_2=\int\psi^4\,dm_2-1>0
\end{equation}
for every $j$. Cauchy--Schwarz gives $\int\psi^4\ge1$. Equality forces
$\psi^2=1$ almost everywhere, hence everywhere by continuity. A real
continuous function valued in $\{-1,1\}$ on a connected surface is constant,
contradicting $\mu>0$. Compactness yields subsequential weak limits of these
probability measures, and \eqref{20003326:obstruction} excludes volume for each
such limit. There is no mass escape and no need to know QUE on the first
factor.

Thus total high energy plus joint Hecke symmetry is not enough when
independent product quotients and a fixed-energy marginal are admitted.
If irreducibility or factorwise high energy is part of the source convention,
this example is excluded. It must not then be advertised as disproving the
catalogue target.

\paragraph{The precise pure-tensor repair.}
Suppose $Y_1,Y_2$ are compact and $v_j,w_j$ are normalized with
\[
 |v_j|^2m_1\rightharpoonup m_1,\qquad
 |w_j|^2m_2\rightharpoonup m_2.
\]
Then $|v_j(x)w_j(y)|^2m_1m_2\rightharpoonup m_1m_2$. For a test
$b(x)c(y)$ the integral is a product of two convergent integrals. Finite sums
of such functions form a uniformly dense algebra on the compact product,
and probability mass one bounds the approximation error independently of
$j$. This proves the assertion by Stone--Weierstrass without an infinite
series interchange. Conversely, convergence of product measures forces
convergence of both marginals by testing functions of a single variable.
Hence marginal equidistribution is necessary and sufficient for these
pure tensors. No comparison of the two energy growth rates is needed once
both marginal conclusions are available.

This proof also shows what remains for joint-eigenspace combinations.
For $u_j=\sum_\alpha c_{j,\alpha}v_{j,\alpha}w_{j,\alpha}$, the density
contains off-diagonal terms
\[
 c_{j,\alpha}\overline{c_{j,\beta}}
 v_{j,\alpha}\overline{v_{j,\beta}}
 w_{j,\alpha}\overline{w_{j,\beta}}.
\]
Marginal diagonal QUE alone does not control their sums, especially if
multiplicity grows. A theorem for every joint eigenfunction needs that
additional information or an argument excluding such multiplicity effects.
No bound on these cross terms is supplied here.

\paragraph{Stress test.}
The example uses compact congruence quotients and genuine joint eigenvectors,
so neither cuspidal normalization nor lack of Hecke symmetry explains the
failure. Its cause is reducibility with a fixed-energy factor. The 2026
proof uses irreducibility at the step unavailable in this construction.
\eref{20003326}{1} On a noncompact product, the probability/tightness argument
would additionally require exclusion of mass escape; it was not silently
extended there. The original catalogue's source-regime qualification is
respected by keeping the result scoped, rather than adopting the broadest
possible interpretation as a purported major counterexample.

\paragraph{Outcome.}
\textbf{Outcome: Exploratory approach with unresolved gap.}
An explicit fixed positive test discrepancy proves failure of the unrestricted
product formulation. A necessary-and-sufficient marginal criterion has also
been proved for pure tensors. These elementary results identify hypotheses
that an exact QUE formulation must retain; no historical novelty is claimed.
The general source-qualified arithmetic QUE target, particularly on
irreducible quotients, is not resolved by this construction.

% END RESEARCH ADDITION 137
\subsection{Sources}
\begin{itemize}\raggedright
\item[\textnormal{[S1]}]\hypertarget{source:20003326:1}{} AIM workshop problem list, Emerging Applications of Measure Rigidity.
\url{https://aimath.org/WWN/measrigid/measrigid.pdf}.
{\small\textit{Catalogue formulation source.}}
% BEGIN RESEARCH REFERENCES 137
\item[\textnormal{[E1]}]\hypertarget{extra:20003326:1}{} Z. Shem-Tov and L. Silberman,
\emph{Arithmetic quantum unique ergodicity for products of hyperbolic 2- and
3-spaces}, Journal d'Analyse Math\'ematique 158 (2026), 413--448;
arXiv:2206.05955, Theorem 1, Sections 2.1 and 3 (quaternionic construction
and the irreducibility step). \url{https://arxiv.org/html/2206.05955}.

% END RESEARCH REFERENCES 137
\end{itemize}

\end{document}
